% Background and Motivation
% Covers: project-based courses, existing tools, research gap
% See .claude/concept.md sections 2 and 7 before writing this chapter.

\label{sec:background}

\subsection{Project-Based Courses and Their Administrative Structure}

Project-based learning engages students in authentic, extended team projects that simulate professional work environments, and is employed across disciplines from software engineering to industrial design, business management, and healthcare education~\cite{pbl}.
These courses share a common administrative structure: students apply or are enrolled, are allocated to project teams, progress through iterative project phases with milestone-based checkpoints, and receive formative feedback throughout before a final summative evaluation.
The iPraktikum practical course at \university{} illustrates this structure:
each semester, over 100 students apply to the course, undergo interviews, are allocated to client-sponsored project teams, and progress through multiple development iterations with assessment at each milestone.
Managing this lifecycle across dozens of teams, multiple tutors, and external clients represents significant administrative overhead that directly detracts from instructional time.

\subsection{Existing Course Administration Tools}

Learning management systems represent the dominant approach to course administration in higher education.
Moodle~\cite{moodle} and Canvas provide content management, gradebook functionality, and plugin ecosystems, but model courses as flat hierarchies of resources and assignments with no concept of administrative lifecycle phases or team-specific workflows.
Within computing education, Artemis~\cite{artemis} provides sophisticated exercise management with automated grading, detailed feedback, and exam workflows, but is designed for individual exercise-based learning and does not address the organizational administration of project teams or course lifecycle phases.
GitHub Classroom~\cite{github-classroom} automates repository creation and assignment distribution for team-based projects but provides no administrative workflow model beyond repository setup.
Ad-hoc combinations of general-purpose tools -- Google Forms for applications, spreadsheets for team allocation, email for feedback -- are widely used but require per-semester manual reconstruction and provide no unified view of the course lifecycle.

Table~\ref{tab:comparison} summarizes the capabilities of these tools against ten administrative requirements derived from project-based course management.

\subsection{Research Gap}

The survey of existing tools reveals a consistent pattern:
tools either address content delivery and individual assessment (LMS, exercise platforms) or handle one specific administrative task (repository creation, team decision support), but none model the course lifecycle as a composable structure.
The result is a gap between the administrative reality of project-based courses -- which unfold as ordered sequences of distinct phases -- and the tooling available to instructors, who must bridge this gap with manual effort every semester.
PROMPT addresses this gap by introducing the phase-based course meta-model as a first-class design abstraction, and providing an open-source platform that operationalizes this model in a production-ready system.
